% Generated draft: requires mathematical and source review.
\documentclass[11pt]{article}
\usepackage[margin=25mm]{geometry}
\usepackage{fontspec}
\setmainfont{Latin Modern Roman}
\usepackage{amsmath,amssymb,mathtools,unicode-math}
\setmathfont{Latin Modern Math}
\usepackage{enumitem,needspace,xurl,hyperref}
\setlength{\parindent}{0pt}
\setlength{\parskip}{5pt}
\setlength{\emergencystretch}{3em}
\newcommand{\N}{\mathbb{N}}
\newcommand{\Z}{\mathbb{Z}}
\newcommand{\Q}{\mathbb{Q}}
\newcommand{\R}{\mathbb{R}}
\newcommand{\C}{\mathbb{C}}
\newcommand{\F}{\mathbb{F}}
\newcommand{\A}{\mathbb{A}}
\newcommand{\PP}{\mathbb{P}}
\newcommand{\E}{\mathbb{E}}
\newcommand{\eps}{\varepsilon}
\newcommand{\dd}{\,\mathrm{d}}
\begin{document}
\section*{UnsolvedMath ID: 700490}
\textbf{Source problem code:} catalogue entry
\subsection{Definitions and mathematical statement}
For a group \(G\), write \(|x|\) for the order of \(x\in G\), and use the commutator convention \([a,b]=a^{-1}b^{-1}ab\). An involution is an element of order exactly \(2\). A conjugacy class of involutions is a set \(D=\{g^{-1}tg:g\in G\}\), where \(|t|=2\).

Describe the finite groups \(G\), up to isomorphism, for which there exists a conjugacy class \(D\) of involutions such that
\[
G=\langle D\rangle
\qquad\text{and}\qquad
\forall a,b\in D,\quad |ab|=4\ \Longrightarrow\ [a,b]\in D.
\]
For involutions, the specified commutator is \([a,b]=(ab)^2\).

The source additionally asserts that this condition holds when \(D\) is a conjugacy class of involutions in a known finite simple group and \(\langle C_D(a)\rangle\) is a \(2\)-group for every \(a\in D\). Here the notation is interpreted as
\[
C_D(a)=\{d\in D:da=ad\},
\]
and a finite \(2\)-group is a group of order \(2^n\) for some integer \(n\geq 0\). This assertion is retained as supplied context; the source record does not specify the scope or date of the word ``known.'' It is not an additional hypothesis restricting the requested classification.
\subsection{Short English statement}
Classify the finite groups generated by a single conjugacy class of involutions such that the commutator of any two members whose product has order \(4\) also belongs to that class. The supplied record labels the problem open; this draft does not independently verify that status.
\subsection{Sources}
\begin{itemize}
\item Local UnsolvedMath registry, supplied record with permanent ID \(700490\), exact source code \texttt{catalogue entry}, titled ``Kourovka Notebook Problem 8.40.'' The supplied status is \texttt{open}; the record was updated on 2026-10-16.
\item \emph{The Kourovka Notebook (Unsolved Problems in Group Theory)}, original problem 8.40, 8th issue (1982); proposed by A. A. Makhn\"{e}v.
\item The supplied background identifies the 21st issue edition (2026), arXiv:1401.0300v48. The collection description names E. I. Khukhro and V. D. Mazurov as editors. These bibliographic details are supplied attribution and have not been independently checked here.
\end{itemize}
\end{document}
